\section{Sequences and Series}

\subsection{Limits of Sequences}

\textbf{Overview.}
This chapter shifts focus from continuous calculus (integrals, derivatives) to \emph{discrete} mathematics: sequences and series.
A sequence assigns a real number to each positive integer; a series asks whether the ``infinite sum'' of those numbers makes sense.
Why care?  Many quantities in science and engineering (power series, Fourier coefficients, numerical approximations) are defined as limits of sequences or sums of series.
Before we can manipulate them, we need a rigorous definition of what ``converges'' means.

\begin{definition}{Sequence and notation}{def:seq-notation}

A (real) \emph{sequence} is a function $a:\mathbb{N}\to\mathbb{R}$, written as $\{a_n\}_{n=1}^{\infty}$ (or simply $\{a_n\}$), where $a_n$ is the $n$-th term.

The starting index does not have to be $1$ (e.g.\ $\{a_n\}_{n=0}^{\infty}$ is also common).

\end{definition}

\noindent
\textbf{Why ``a function from $\mathbb{N}$ to $\mathbb{R}$''?}
At first glance, a sequence looks like a list: $a_1,a_2,a_3,\dots$
But mathematically, a list \emph{is} a function: you input the position $n$ and output the value $a_n$.
Thinking of a sequence as a function is not pedantic; it lets us immediately apply function-theoretic tools (e.g.\ L'H\^opital's Rule via Theorem~\ref{thm:seq-by-function}) to compute sequence limits.

\paragraph{Geometric intuition.}

Think of a sequence as a list of points on the number line (or as plotted points $(n,a_n)$).

Saying ``$a_n\to L$'' means that, as $n$ grows, the points $a_n$ are eventually trapped inside \emph{any} narrow band around $L$.

% Figure description (fallback): Plot the terms a_n = n/(n+1) as points (n,a_n) approaching the horizontal line y=1.

\begin{center}
\begin{tikzpicture}
\begin{axis}[
    width=0.78\linewidth,
    height=5cm,
    xmin=0, xmax=11,
    ymin=0.4, ymax=1.1,
    axis lines=left,
    xlabel={$n$},
    ylabel={$a_n$},
    ytick={0.5,0.6,0.7,0.8,0.9,1.0},
    xtick={0,2,4,6,8,10},
    ticklabel style={font=\small},
    label style={font=\small},
] % <--- This bracket was missing!

% The sequence points (n, a_n)
\addplot[
    only marks, 
    mark=*, 
    blue, 
    samples at={1,2,...,10}
] {x/(x+1)};

% The horizontal asymptote at y=1
\addplot[
    dashed, 
    red, 
    domain=0:11
] {1} node[right, pos=1] {$y=1$};

\end{axis}
\end{tikzpicture}
\end{center}

\begin{definition}{Limit of a sequence (formal $\varepsilon$--$N$ definition)}{def:seq-limit}

We say $\displaystyle \lim_{n\to\infty} a_n = L$ (or $a_n\to L$) if for every $\varepsilon>0$ there exists an integer $N$ such that

\[
n>N \ \Longrightarrow\ |a_n-L|<\varepsilon.
\]

If such an $L$ exists, $\{a_n\}$ is \emph{convergent}; otherwise it is \emph{divergent}.

\end{definition}

\paragraph{Why this definition looks the way it does.}

The quantity $|a_n-L|$ is the distance from $a_n$ to $L$. The phrase ``for every $\varepsilon>0$'' means:

no matter how strict the tolerance is, the tail of the sequence eventually stays within that tolerance.

The index $N$ may depend on $\varepsilon$ (tighter tolerance $\Rightarrow$ larger $N$).

\paragraph{Unpacking the definition further (exam-ready).}
The definition has three quantifiers, and their \emph{order} matters:
\begin{enumerate}[leftmargin=3em]
\item \textbf{``For every $\varepsilon>0$'':} the challenger picks any positive tolerance (no matter how small).
\item \textbf{``There exists $N$'':} you must respond with a cutoff index $N$ (which may depend on $\varepsilon$).
\item \textbf{``For all $n>N$'':} every term past the cutoff must satisfy $|a_n-L|<\varepsilon$.
\end{enumerate}
Think of it as a game: if you can always win (i.e.\ always find an $N$), the sequence converges to $L$.

\begin{examtip}[title={Exam Focus: The $\varepsilon$--$N$ proof template}]

For an $\varepsilon$--$N$ proof, always follow this structure:

\begin{enumerate}[leftmargin=3em]
\item State ``Let $\varepsilon>0$.''
\item Compute $|a_n-L|$ and simplify it.
\item Find an upper bound for $|a_n-L|$ that is a simple decreasing function of $n$ (e.g.\ $\frac{C}{n}$ or $\frac{C}{n^2}$).
\item Choose $N$ so that the upper bound $<\varepsilon$ whenever $n\ge N$ (e.g.\ $N=\lceil C/\varepsilon\rceil$).
\item Conclude: ``For all $n\ge N$, $|a_n-L|\le \text{bound}\le \varepsilon$.''
\end{enumerate}

The ``backward work'' (finding the bound and choosing $N$) is \emph{scratch work} that does not appear in the formal proof; the formal proof runs \emph{forward} from the chosen $N$.

\end{examtip}


\begin{theorem}{Uniqueness and boundedness}{uniq-bounded}
\begin{enumerate}[label=(\roman*), leftmargin=3em]
\item A convergent sequence has a unique limit.

\item Every convergent sequence is bounded.

\end{enumerate}

\end{theorem}

\noindent
\textbf{Why uniqueness?}
If $a_n\to L$ and $a_n\to L'$, then by the triangle inequality $|L-L'|\le |a_n-L|+|a_n-L'|<2\varepsilon$ for $n$ large enough.
Since $\varepsilon$ was arbitrary, $|L-L'|=0$, so $L=L'$.

\noindent
\textbf{Why boundedness?}
If $a_n\to L$, take $\varepsilon=1$: there exists $N$ such that $|a_n-L|<1$ for all $n>N$, i.e.\ $a_n\in(L-1,L+1)$.
The finitely many terms $a_1,\dots,a_N$ are just finitely many real numbers, so they are bounded.
Hence the \emph{entire} sequence is bounded.

\begin{commonmistake}[title={Common Mistake: Converse of boundedness}]

Theorem~\ref{thm:uniq-bounded}(ii) says convergent $\Rightarrow$ bounded, but the converse is \textbf{false}.

Example: $a_n=(-1)^n$ is bounded ($|a_n|\le 1$) but divergent (oscillates between $1$ and $-1$).

You need \emph{monotonicity + boundedness} (Theorem~\ref{thm:monotone}) to guarantee convergence.

\end{commonmistake}

\begin{examplebox}[title={Worked Example (Tutorial 7, Question 1(a)): $\varepsilon$--$N$ proof for $a_n=\frac{1}{3n}$}]

Show that $\displaystyle \lim_{n\to\infty}\frac{1}{3n}=0$.

\emph{Solution.}

Let $\varepsilon>0$. We want $|a_n-0|<\varepsilon$, i.e.

\[
\left|\frac{1}{3n}\right|<\varepsilon
\quad\Longleftrightarrow\quad
n>\frac{1}{3\varepsilon}.
\]

Choose $N=\left\lceil \frac{1}{3\varepsilon}\right\rceil$. Then for all $n\ge N$,

\[
|a_n-0|=\frac{1}{3n}\le \frac{1}{3N}<\varepsilon.
\]

Hence $\displaystyle \boxed{\lim_{n\to\infty}\frac{1}{3n}=0}$.

\end{examplebox}

\begin{examplebox}[title={Worked Example (Tutorial 7, Question 1(c)): $\varepsilon$--$N$ proof for a rational sequence}]

Show that $\displaystyle \lim_{n\to\infty}\frac{2n^2-3n}{3n^2+1}=\frac{2}{3}$.

\emph{Solution.}

Let $\varepsilon>0$. Compute the difference:

\[
\left|\frac{2n^2-3n}{3n^2+1}-\frac{2}{3}\right|
=\left|\frac{3(2n^2-3n)-2(3n^2+1)}{3(3n^2+1)}\right|
=\frac{9n+2}{9n^2+3}.
\]

\textbf{Bounding step.} For $n\ge 1$:

\[
\frac{9n+2}{9n^2+3}\le \frac{9n+2}{9n^2}
=\frac{1}{n}+\frac{2}{9n^2}
\le \frac{1}{n}+\frac{2}{9n}
=\frac{11}{9n}.
\]

Choose $N=\left\lceil \frac{11}{9\varepsilon}\right\rceil$. Then for all $n\ge N$:

\[
\left|\frac{2n^2-3n}{3n^2+1}-\frac{2}{3}\right|\le \frac{11}{9n}\le \frac{11}{9N}<\varepsilon.
\]

Hence $\displaystyle \boxed{\lim_{n\to\infty}\frac{2n^2-3n}{3n^2+1}=\frac{2}{3}}$.

\textbf{Exam note:} The key step is replacing the exact expression $\frac{9n+2}{9n^2+3}$ with a simpler upper bound $\frac{C}{n}$. Any valid bound works; you do \emph{not} need to find the tightest one.

\end{examplebox}

\begin{definition}{Subsequence}{def:subsequence}

Given $\{a_n\}_{n=1}^{\infty}$, a \emph{subsequence} is a sequence of the form

\[
a_{n_1},a_{n_2},a_{n_3},\dots
\quad\text{where}\quad
1\le n_1<n_2<n_3<\cdots.
\]

\end{definition}

\noindent
\textbf{Why subsequences matter.}
The subsequence test (below) is the workhorse for \emph{proving divergence}: if you can exhibit two subsequences with different limits, the original sequence cannot converge.
This is typically much easier than a direct $\varepsilon$--$N$ argument for divergence.

\begin{theorem}{Subsequence test}{thm:subseq}

If $\{a_n\}$ converges to $L$, then \emph{every} subsequence of $\{a_n\}$ also converges to $L$.

\end{theorem}

\noindent
\textbf{Why this is true (proof sketch).}
If $n>N$ implies $|a_n-L|<\varepsilon$, and $\{n_k\}$ is a strictly increasing sequence of indices, then $n_k\ge k$ for all $k$.
So for $k>N$, we have $n_k\ge k>N$, giving $|a_{n_k}-L|<\varepsilon$.

\noindent
\textbf{Contrapositive (the useful direction for exams):}

\begin{itemize}[leftmargin=*]
\item If two subsequences converge to \emph{different} limits $\Rightarrow$ $\{a_n\}$ diverges.
\item If any one subsequence diverges $\Rightarrow$ $\{a_n\}$ diverges.
\end{itemize}

\begin{examplebox}[title={Worked Example: $\{(-1)^n\}$ diverges (Lecture Example 4.4)}]

The even subsequence $a_{2k}=(-1)^{2k}=1\to 1$.

The odd subsequence $a_{2k-1}=(-1)^{2k-1}=-1\to -1$.

Since $1\ne -1$, by the subsequence test, $\{(-1)^n\}$ diverges.

\end{examplebox}

\begin{examplebox}[title={Worked Example (Tutorial 7, Question 2(vii)): subsequence test for $\frac{(-1)^n n^3}{n^3+2n^2+1}$}]

Write $a_n=(-1)^n\cdot \dfrac{1}{1+\frac{2}{n}+\frac{1}{n^3}}$.

\begin{itemize}[leftmargin=*]
\item Even subsequence: $a_{2k}=+\dfrac{1}{1+\frac{2}{2k}+\frac{1}{(2k)^3}}\to 1$.
\item Odd subsequence: $a_{2k-1}=-\dfrac{1}{1+\frac{2}{2k-1}+\frac{1}{(2k-1)^3}}\to -1$.
\end{itemize}

Since $1\ne -1$, the sequence diverges by the subsequence test.

\end{examplebox}

\begin{examplebox}[title={Worked Example (Tutorial 7, Q3): If odd and even subsequences both converge to $L$, then $a_n\to L$}]

\textbf{Claim.} If $a_{2k-1}\to L$ and $a_{2k}\to L$, then $a_n\to L$.

\emph{Proof.}
Let $\varepsilon>0$. Since $a_{2k}\to L$, there exists $N_e$ such that $|a_{2k}-L|<\varepsilon$ for all $k\ge N_e$. Similarly, there exists $N_o$ such that $|a_{2k-1}-L|<\varepsilon$ for all $k\ge N_o$.

Let $N=\max\{2N_e,\,2N_o-1\}$. For any $n\ge N$:
\begin{itemize}[leftmargin=*]
\item If $n=2k$ (even): $n\ge 2N_e$ implies $k\ge N_e$, so $|a_n-L|<\varepsilon$.
\item If $n=2k-1$ (odd): $n\ge 2N_o-1$ implies $k\ge N_o$, so $|a_n-L|<\varepsilon$.
\end{itemize}
Hence $a_n\to L$.\ $\square$

\textbf{Exam use:} This fact is useful for sequences defined piecewise (e.g.\ even and odd terms given by different formulas). If both halves converge to the same $L$, the full sequence converges.

\end{examplebox}

\begin{definition}{Divergence to $\pm\infty$}{def:div-infty}

We say $a_n\to\infty$ if for every $M>0$ there exists $N$ such that $n>N\Rightarrow a_n>M$.

Similarly define $a_n\to -\infty$.

\end{definition}

\noindent
\textbf{Why a separate definition?}
The $\varepsilon$--$N$ definition of $a_n\to L$ requires $L$ to be a \emph{real number}. Saying $a_n\to\infty$ is \emph{not} the same as saying ``$a_n$ converges to $\infty$''---it is a specific type of divergence where the terms grow without bound.

\begin{keyformula}[title={Quick implications for positive sequences}]

Assume $a_n>0$ for all sufficiently large $n$.

\begin{enumerate}[label=(\roman*), leftmargin=3em]

\item If $a_n\to 0$, then $\displaystyle \frac{1}{a_n}\to \infty$.

\item If $a_n\to \infty$, then $\displaystyle \frac{1}{a_n}\to 0$.

\item If $b_n\ge a_n$ eventually and $a_n\to\infty$, then $b_n\to\infty$.

\end{enumerate}

\end{keyformula}

\noindent
\textbf{Why these hold (intuition).}
(i) If $a_n$ is very small and positive, $1/a_n$ is very large.
(ii) If $a_n$ is very large, $1/a_n$ is very small.
(iii) If $a_n$ already exceeds every bound, and $b_n\ge a_n$, then $b_n$ does too.

\begin{commonmistake}

Divergence does \emph{not} behave well under algebra.

For instance, $a_n=(-1)^n$ is divergent, but $|a_n|=1$ is convergent.

Always justify convergence before applying limit laws.

\end{commonmistake}

\begin{examtip}

To prove divergence quickly, the subsequence test is often the fastest:

\begin{itemize}

\item Try even/odd subsequences when $(-1)^n$, $\cos(n\pi)$, or oscillations appear.

\item If the expression behaves differently along two simple subsequences (e.g.\ $n=2k$ and $n=2k+1$), you can often conclude divergence immediately.

\end{itemize}

\end{examtip}

%--------------------------------------------------

\subsection{Finding Limits of Sequences}

\textbf{Motivation.}
The $\varepsilon$--$N$ definition tells us what convergence \emph{means}, but using it directly for every limit would be tedious. This section builds a toolkit of shortcuts: limit laws, the squeeze theorem, the function--sequence bridge (Theorem~\ref{thm:seq-by-function}), and the continuity transfer (Theorem~\ref{thm:cont-seq}). Together, these let you compute most sequence limits algebraically.

\begin{keyformula}[title={Limit laws for sequences (require convergence)}]

Suppose $\{a_n\}$ and $\{b_n\}$ are convergent and $c\in\mathbb{R}$ is a constant. Then:

\begin{enumerate}[label=(\alph*), leftmargin=3em]

\item $\lim (a_n\pm b_n)=\lim a_n \pm \lim b_n$.

\item $\lim (c a_n)=c\lim a_n$.

\item $\lim (a_n b_n)=(\lim a_n)(\lim b_n)$.

\item $\lim \dfrac{a_n}{b_n}=\dfrac{\lim a_n}{\lim b_n}$, provided $\lim b_n\ne 0$.

\end{enumerate}

\end{keyformula}

\noindent
\textbf{Why the convergence hypothesis is non-negotiable.}
Without it, absurd conclusions follow. For example, $a_n=n$ and $b_n=-n$ are both divergent, but $a_n+b_n=0\to 0$. Writing ``$\lim(a_n+b_n)=\lim a_n+\lim b_n=\infty+(-\infty)$'' is meaningless.

\paragraph{Method selection (high yield).}

When asked for $\lim_{n\to\infty} a_n$, try in this order:

\begin{enumerate}[leftmargin=3em]

\item \textbf{Algebraic simplification}: factor/divide by the dominant power of $n$; rationalize radicals.

\item \textbf{Squeeze}: use boundedness ($|\sin|,|\cos|\le 1$), or compare with simpler sequences.

\item \textbf{Recognise a function}: if $a_n=f(n)$ and $\lim_{x\to\infty} f(x)$ is known, transfer the result.

\item \textbf{Continuity}: if $a_n\to L$, then $f(a_n)\to f(L)$ for continuous $f$.

\end{enumerate}

\begin{examtip}[title={Exam Focus: When to use which method}]

\begin{itemize}[leftmargin=*]
\item \textbf{Rational expressions in $n$:} Divide top and bottom by $n^{\text{highest power}}$. This always works and is the fastest approach.
\item \textbf{Bounded oscillating factor $\times$ vanishing term:} Squeeze (e.g.\ $\frac{\sin n}{n}$, $\frac{\cos^2 n}{2n}$).
\item \textbf{Expression involving $\ln$, $e$, $n!$:} Recognise as $f(n)$ and use L'H\^opital's Rule on the continuous version $f(x)$.
\item \textbf{Composition with continuous function:} Use Theorem~\ref{thm:cont-seq} (e.g.\ $\sin(\pi/n)$, $e^{3n/(n+1)}$).
\end{itemize}
\end{examtip}

\begin{theorem}{Sequence defined by a function}{thm:seq-by-function}

Let $f:(A,\infty)\to\mathbb{R}$ be a function such that $\lim_{x\to\infty} f(x)=L$, and define $a_n=f(n)$ for integers $n$ large enough.

Then $\lim_{n\to\infty} a_n = L$.

\end{theorem}

\noindent
\textbf{Why this works.}
The real limit $\lim_{x\to\infty} f(x)=L$ means: for every $\varepsilon>0$, there exists $M$ such that $x>M\Rightarrow |f(x)-L|<\varepsilon$.
Integers are a subset of the reals, so choosing $N=\lceil M\rceil$ gives $|f(n)-L|<\varepsilon$ for all $n\ge N$.

\textbf{Practical payoff:} This lets you apply L'H\^opital's Rule (which requires differentiability, a function concept) to compute sequence limits.

\begin{theorem}{Squeeze theorem for sequences}{thm:squeeze-seq}

If $a_n\le b_n\le c_n$ for all $n\ge N_0$ and $\lim a_n=\lim c_n=L$, then $\lim b_n=L$.

\end{theorem}

\noindent
\textbf{Why the squeeze works.}
Given $\varepsilon>0$, for $n$ large enough, $L-\varepsilon<a_n\le b_n\le c_n<L+\varepsilon$, so $|b_n-L|<\varepsilon$.
The inequalities $a_n\le b_n\le c_n$ ``squeeze'' $b_n$ into the $\varepsilon$-band around $L$.

\begin{theorem}{Continuity and limits of sequences}{thm:cont-seq}

If $a_n\to L$ and $f$ is continuous at $L$, then $f(a_n)\to f(L)$.

\end{theorem}

\noindent
\textbf{Why this is true.}
Continuity at $L$ means: for every $\varepsilon>0$, there exists $\delta>0$ such that $|x-L|<\delta\Rightarrow |f(x)-f(L)|<\varepsilon$.
Since $a_n\to L$, eventually $|a_n-L|<\delta$, so $|f(a_n)-f(L)|<\varepsilon$.

\begin{examplebox}[title={Worked Example: geometric sequence $r^n$}]

Determine $\displaystyle \lim_{n\to\infty} r^n$.

\emph{Solution.}

\begin{enumerate}[label=(\roman*), leftmargin=3em]

\item If $|r|<1$, then $|r|^n\to 0$, so $r^n\to 0$.

\item If $r=1$, then $r^n=1$ for all $n$.

\item If $r>1$, then $r^n\to\infty$ (unbounded growth).

\item If $r\le -1$, the terms do not approach a single real number (they oscillate in sign and/or grow in magnitude), hence divergence.

\end{enumerate}

\end{examplebox}

\noindent
\textbf{Why the case $|r|<1$ works.}
If $0<r<1$, write $r=\frac{1}{1+p}$ with $p>0$. By Bernoulli's inequality, $(1+p)^n\ge 1+np$, so $0<r^n=\frac{1}{(1+p)^n}\le \frac{1}{1+np}\to 0$.
For $-1<r<0$, we have $|r^n|=|r|^n\to 0$, and $|a_n|\to 0\Rightarrow a_n\to 0$ by the squeeze theorem.

\begin{examplebox}[title={Worked Example: a bounded trigonometric factor (Squeeze)}]

Evaluate $\displaystyle \lim_{n\to\infty}\frac{\cos^2 n}{2n}$.

\emph{Solution.}

Since $0\le \cos^2 n\le 1$,

\[
0\le \frac{\cos^2 n}{2n}\le \frac{1}{2n}.
\]

As $n\to\infty$, $\frac{1}{2n}\to 0$. Hence by Theorem~\ref{thm:squeeze-seq},

\[
\boxed{\lim_{n\to\infty}\frac{\cos^2 n}{2n}=0}.
\]

\end{examplebox}

\begin{examplebox}[title={Worked Example: continuity trick}]

Evaluate $\displaystyle \lim_{n\to\infty}\sin\!\left(\frac{\pi}{n}\right)$.

\emph{Solution.}

We know $\frac{\pi}{n}\to 0$. Since $\sin x$ is continuous at $0$, Theorem~\ref{thm:cont-seq} gives

\[
\boxed{\lim_{n\to\infty}\sin\!\left(\frac{\pi}{n}\right)=\sin 0=0.}
\]

\end{examplebox}

\begin{examplebox}[title={Worked Example (Tutorial 7, Q2(viii)): continuity with $\tan$}]

Evaluate $\displaystyle \lim_{n\to\infty}\tan\!\left(\frac{2n\pi}{1+8n}\right)$.

\emph{Solution.}
First find the ``inner'' limit:
\[
\frac{2n\pi}{1+8n}=\frac{2\pi}{8+\frac{1}{n}}\to \frac{2\pi}{8}=\frac{\pi}{4}.
\]
Since $\tan$ is continuous at $\pi/4$ (which is not an odd multiple of $\pi/2$), Theorem~\ref{thm:cont-seq} gives
\[
\lim_{n\to\infty}\tan\!\left(\frac{2n\pi}{1+8n}\right)=\tan\!\left(\frac{\pi}{4}\right)=\boxed{1}.
\]

\textbf{Sanity check:} Always verify the inner limit lands at a point where the outer function is continuous. If it landed at $\pi/2$, $\tan$ would be undefined and the argument would fail.

\end{examplebox}

\begin{examplebox}[title={Worked Example: comparing growth rates}]

Show that $\displaystyle \lim_{n\to\infty}\frac{n!}{n^n}=0$.

\emph{Solution.}

\[
\frac{n!}{n^n}=\frac{1\cdot 2\cdot 3\cdots n}{n\cdot n\cdot n\cdots n}
=\prod_{k=1}^{n}\frac{k}{n}.
\]
For $2\le k\le n$, $\frac{k}{n}\le 1$, so
\[
0\le \frac{n!}{n^n}=\frac{1}{n}\cdot\prod_{k=2}^{n}\frac{k}{n}\le \frac{1}{n}.
\]

Since $\frac{1}{n}\to 0$, the squeeze theorem yields $\displaystyle \boxed{\lim_{n\to\infty}\frac{n!}{n^n}=0}$.

\end{examplebox}

\begin{examplebox}[title={Worked Example: if $|a_n|\to 0$ then $a_n\to 0$}]

Assume $\lim_{n\to\infty}|a_n|=0$. Show $\lim_{n\to\infty}a_n=0$.

\emph{Solution.}

For every $n$, $-|a_n|\le a_n\le |a_n|$.

Taking limits, the squeeze theorem gives $a_n\to 0$.

\end{examplebox}

\begin{examplebox}[title={Worked Example (Lecture Example 4.14): L'H\^opital via the function bridge}]

Calculate $\displaystyle \lim_{n\to\infty}\frac{n+\ln n}{n^2}$.

\emph{Solution.}
Define $f(x)=\frac{x+\ln x}{x^2}$. This is an $\frac{\infty}{\infty}$ form as $x\to\infty$. By L'H\^opital's Rule:
\[
\lim_{x\to\infty}\frac{x+\ln x}{x^2}
=\lim_{x\to\infty}\frac{1+\frac{1}{x}}{2x}
=\lim_{x\to\infty}\frac{1}{2x}+\frac{1}{2x^2}=0.
\]
By Theorem~\ref{thm:seq-by-function}, $\boxed{\lim_{n\to\infty}\frac{n+\ln n}{n^2}=0}$.

\end{examplebox}

\begin{examplebox}[title={Worked Example (Lecture Example 4.15): composing function bridge with continuity}]

Calculate $\displaystyle \lim_{n\to\infty}e^{3n/(n+1)}$.

\emph{Solution.}
Let $a_n=\frac{3n}{n+1}=\frac{3}{1+1/n}\to 3$.

Since $f(x)=e^x$ is continuous at $3$, Theorem~\ref{thm:cont-seq} gives
\[
\lim_{n\to\infty}e^{3n/(n+1)}=e^{\lim a_n}=\boxed{e^3}.
\]

\end{examplebox}

\begin{examplebox}[title={Worked Example (Tutorial 7, Q5): $\sqrt[n]{a}\to 1$ for $a>0$}]

\textbf{Claim:} $\displaystyle \lim_{n\to\infty}a^{1/n}=1$ for all $a>0$.

\emph{Proof (continuity method).}
Write $a^{1/n}=e^{\frac{\ln a}{n}}$. Since $\frac{\ln a}{n}\to 0$ and $e^x$ is continuous at $0$:
\[
\lim_{n\to\infty}a^{1/n}=e^0=1.
\]

\emph{Proof (squeeze for $a>1$).}
Let $b=a-1>0$. By Bernoulli's inequality, $(1+b/n)^n\ge 1+b=a$, so $a^{1/n}\le 1+b/n$.
Also $a^{1/n}\ge 1$ since $a>1$. Hence $1\le a^{1/n}\le 1+\frac{b}{n}\to 1$ by squeeze.

For $0<a<1$, apply the $a>1$ case to $1/a$ and take reciprocals.

\end{examplebox}

\begin{extension}[title={Extension (Beyond Syllabus: Stolz--Ces\`aro theorem for limits)}]

\textbf{Syllabus link:} Chapter 4 --- \S4.2 Finding Limits of Sequences.

\medskip

\textbf{Theorem (Stolz--Ces\`aro).}

Let $\{b_n\}$ be a strictly increasing sequence with $b_n\to\infty$.

Let $\{a_n\}$ be any sequence.

If the limit

\[
\lim_{n\to\infty}\frac{a_{n+1}-a_n}{b_{n+1}-b_n}=L
\]

exists (finite), then $\displaystyle \lim_{n\to\infty}\frac{a_n}{b_n}=L$.

\medskip

\textbf{Application.} Compute $\displaystyle \lim_{n\to\infty}\frac{\ln n}{n}$.

Take $a_n=\ln n$ and $b_n=n$. Then $b_n$ is strictly increasing and $b_n\to\infty$.

Moreover,

\[
\frac{a_{n+1}-a_n}{b_{n+1}-b_n}
=\frac{\ln(n+1)-\ln n}{(n+1)-n}
=\ln\!\left(1+\frac{1}{n}\right).
\]
Since $\ln(1+x)\to 0$ as $x\to 0$, we have $\ln\!\left(1+\frac{1}{n}\right)\to 0$.
By Stolz--Ces\`aro,
\[
\boxed{\lim_{n\to\infty}\frac{\ln n}{n}=0.}
\]

\end{extension}

\begin{commonmistake}

The limit laws require the individual limits to exist.

If $\{a_n\}$ or $\{b_n\}$ is divergent, statements like ``$\lim(a_n+b_n)=\lim a_n+\lim b_n$'' are not valid.

\end{commonmistake}

\begin{examplebox}[title={Worked Example (Tutorial 7, Q4): True/false on divergent sequences}]

\textbf{(a)} ``If $\{a_n\}$ and $\{b_n\}$ are divergent, then $\{a_n+b_n\}$ is divergent.'' \textbf{FALSE.}

\emph{Counterexample:} $a_n=(-1)^n$, $b_n=-(-1)^n$. Both diverge, but $a_n+b_n=0\to 0$.

\textbf{(b)} ``If $\{a_n\}$ is divergent, then $\{|a_n|\}$ is divergent.'' \textbf{FALSE.}

\emph{Counterexample:} $a_n=(-1)^n$ diverges, but $|a_n|=1\to 1$.

\textbf{(c)} ``If $a_n\to L$ and $b_n\to 0$, then $a_nb_n\to 0$.'' \textbf{TRUE.}

\emph{Proof:} Since $a_n\to L$, $\{a_n\}$ is bounded: $|a_n|\le M$. Given $\varepsilon>0$, choose $N$ so that $|b_n|<\varepsilon/M$ for $n\ge N$. Then $|a_nb_n|\le M|b_n|<\varepsilon$.

\end{examplebox}

\begin{examtip}

For rational expressions in $n$:

\begin{itemize}

\item Divide numerator and denominator by the highest power of $n$ present.

\item For $\sqrt{\text{polynomial in }n}$, factor out the dominant power inside the root: $\sqrt{n^2(1+\cdots)}=n\sqrt{1+\cdots}$.

\item If you rationalize, rationalize \emph{once} and simplify before taking limits.

\end{itemize}

\end{examtip}

%--------------------------------------------------

\subsection{Monotonic Sequences}

\textbf{Motivation.}
Many important sequences (especially recursively defined ones) are hard to evaluate by direct algebraic manipulation.
The Monotone Sequence Theorem gives an \emph{existence} guarantee: if a sequence is monotonic and bounded, it converges---even if you cannot compute the limit directly.
You then find the limit by solving the fixed-point equation $L=f(L)$.
This three-step workflow (bound $\to$ monotone $\to$ fixed point) is one of the most frequently examined techniques in MH1101.

\begin{definition}{Monotonicity}{def:monotone}

A sequence $\{a_n\}$ is \emph{increasing} if $a_n\le a_{n+1}$ for all $n$ (and \emph{strictly increasing} if $a_n<a_{n+1}$).

A sequence $\{a_n\}$ is \emph{decreasing} if $a_n\ge a_{n+1}$ for all $n$ (and \emph{strictly decreasing} if $a_n>a_{n+1}$).

A sequence is \emph{monotonic} if it is increasing or decreasing.

\end{definition}

\noindent
\textbf{Ways to check monotonicity.}
\begin{itemize}[leftmargin=*]
\item \textbf{Difference test:} Compute $a_{n+1}-a_n$ and determine its sign.
\item \textbf{Ratio test (for positive sequences):} Compute $a_{n+1}/a_n$ and check whether it is $\ge 1$ or $\le 1$.
\item \textbf{Function derivative:} If $a_n=f(n)$, compute $f'(x)$ and check its sign.
\item \textbf{For recursive sequences $a_{n+1}=f(a_n)$:} Show $f(x)\ge x$ (or $\le x$) on the relevant interval.
\end{itemize}

\begin{definition}{Boundedness}{def:bounded}

A sequence $\{a_n\}$ is \emph{bounded above} if there exists $M$ such that $a_n\le M$ for all $n$.

It is \emph{bounded below} if there exists $m$ such that $m\le a_n$ for all $n$.

It is \emph{bounded} if it is both bounded above and bounded below.

\end{definition}

\begin{theorem}{Monotone sequence theorem}{thm:monotone}

Every bounded, monotonic sequence is convergent.

\end{theorem}

\noindent
\textbf{Why this theorem is true (intuition).}
An increasing sequence that is bounded above has a least upper bound (supremum) $L$ by the completeness of $\mathbb{R}$. Since $L$ is the \emph{least} upper bound, for any $\varepsilon>0$, there must be some $a_N>L-\varepsilon$ (otherwise $L-\varepsilon$ would be a smaller upper bound). Since the sequence is increasing, all subsequent terms satisfy $a_n\ge a_N>L-\varepsilon$, and also $a_n\le L<L+\varepsilon$. Hence $|a_n-L|<\varepsilon$ for all $n\ge N$.

\paragraph{Standard recursion workflow.}

For a recursively defined sequence $a_{n+1}=f(a_n)$:

\begin{enumerate}[leftmargin=3em]

\item \textbf{Invariant interval (boundedness):} find an interval $I=[m,M]$ such that $a_1\in I$ and $f(I)\subseteq I$.

\item \textbf{Monotonicity:} show $a_{n+1}\ge a_n$ (or $\le$) using algebra and the bounds from Step 1.

\item \textbf{Limit identification:} conclude convergence by Theorem~\ref{thm:monotone}; let $a_n\to L$ and solve $L=f(L)$.

\item \textbf{Reject extraneous roots:} keep only solutions consistent with the bounds (e.g.\ $L\in I$).

\end{enumerate}

\begin{examtip}[title={Exam Focus: Why you \emph{must} prove convergence before solving $L=f(L)$}]

The equation $L=f(L)$ may have multiple solutions, and without knowing the sequence converges, you cannot even write $\lim a_{n+1}=f(\lim a_n)$.
Example: define $a_1=0$, $a_{n+1}=a_n^2+a_n$. The equation $L=L^2+L$ gives $L=0$, but the sequence is $0,0,0,\dots$ \emph{and} $L=0$ is the only fixed point---so it works. But if you had $a_1=2$, the sequence diverges to $\infty$, and ``solving $L=f(L)$'' gives the wrong answer.

Always establish convergence \emph{first}, then solve the fixed-point equation.

\end{examtip}

\begin{examplebox}[title={Worked Example (Tutorial 8, Question 1): nested radicals}]

Find the limit of

\[
\sqrt{2},\ \sqrt{2\sqrt{2}},\ \sqrt{2\sqrt{2\sqrt{2}}},\ \ldots
\]

Define $a_1=\sqrt{2}$ and $a_{n+1}=\sqrt{2a_n}$.

\emph{Bounded:} show by induction that $0<a_n<2$ for all $n$.

Base: $a_1=\sqrt{2}<2$.\ $\checkmark$

Inductive step: if $a_n<2$, then $a_{n+1}=\sqrt{2a_n}<\sqrt{2\cdot 2}=2$.\ $\checkmark$

Also $a_n>0$ trivially (positive under square root).

\emph{Increasing:} $a_{n+1}\ge a_n \iff \sqrt{2a_n}\ge a_n \iff 2a_n\ge a_n^2 \iff a_n(2-a_n)\ge 0$.

Since $0<a_n<2$, both factors are positive. So increasing.\ $\checkmark$

\emph{Limit:} by Theorem~\ref{thm:monotone}, $a_n\to L$ exists. Then $L=\sqrt{2L}$, so $L^2=2L$, giving $L(L-2)=0$.

Since $L>0$, we conclude $\boxed{L=2}$.

\end{examplebox}

\begin{examplebox}[title={Worked Example (Tutorial 8, Question 2): $a_1=1$, $a_{n+1}=3-\frac{1}{a_n}$}]

Show the sequence is increasing and $0<a_n<3$, then find the limit.

\emph{Invariant interval:} We show $a_n\in[1,3)$ by induction.

Base: $a_1=1\in[1,3)$.\ $\checkmark$

Inductive step: if $1\le a_n<3$, then $\frac{1}{a_n}\le 1$, so $a_{n+1}=3-\frac{1}{a_n}\ge 2\ge 1$.

Also $\frac{1}{a_n}>0$, so $a_{n+1}<3$.\ $\checkmark$

\emph{Increasing:} $a_{n+1}-a_n=3-\frac{1}{a_n}-a_n=\frac{-a_n^2+3a_n-1}{a_n}$.

The numerator $-a_n^2+3a_n-1\ge 0$ when $a_n\in[\alpha,\beta]$ where $\alpha=\frac{3-\sqrt{5}}{2}\approx 0.382$ and $\beta=\frac{3+\sqrt{5}}{2}\approx 2.618$.

Since $a_n\in[1,3)\subset[\alpha,\beta]$ (check: $1>\alpha$ and $3>\beta$ but $a_n<3$ and is increasing with $a_n\le\beta$), the numerator is $\ge 0$ and $a_n>0$, so $a_{n+1}\ge a_n$.\ $\checkmark$

\emph{Limit:} $L=3-\frac{1}{L}\Rightarrow L^2-3L+1=0\Rightarrow L=\frac{3\pm\sqrt{5}}{2}$.

Since $a_n$ is increasing from $a_1=1$ and bounded above by $\beta=\frac{3+\sqrt{5}}{2}$:

$\boxed{L=\frac{3+\sqrt{5}}{2}}$.

\end{examplebox}

\begin{examplebox}[title={Worked Example (Tutorial 8, Question 3): $a_1=2$, $a_{n+1}=\frac{1}{3-a_n}$}]

\emph{Invariant interval:} Show $a_n\in(\alpha,\beta]$ where $\alpha=\frac{3-\sqrt{5}}{2}$ and $\beta=\frac{3+\sqrt{5}}{2}$. (Both are fixed points of $f(x)=\frac{1}{3-x}$.)

Since $a_1=2\in(\alpha,\beta)$ and $f$ is increasing on $(-\infty,3)$ with $f(\alpha)=\alpha$, $f(\beta)=\beta$, the interval $[\alpha,\beta]$ is invariant.

\emph{Decreasing:} $a_{n+1}\le a_n\iff \frac{1}{3-a_n}\le a_n\iff a_n^2-3a_n+1\le 0$, which holds for $a_n\in[\alpha,\beta]$.\ $\checkmark$

Since $\{a_n\}$ is decreasing and bounded below, it converges. Solving $L=\frac{1}{3-L}$ gives $L\in\{\alpha,\beta\}$. Since the sequence decreases from $a_1=2$ and $a_2=1<2$, we have $L\le 1<\beta$, so:

$\boxed{L=\frac{3-\sqrt{5}}{2}}$.

\end{examplebox}

\begin{commonmistake}

For recursive sequences, solving $L=f(L)$ is \emph{not} enough.

You must first justify that $\{a_n\}$ converges (typically by boundedness + monotonicity).

\end{commonmistake}

\begin{examtip}

When proving monotonicity, avoid expanding blindly.

A reliable pattern is:

\[
a_{n+1}\ge a_n \ \Longleftrightarrow\ f(a_n)\ge a_n,
\]

then factor and use bounds on $a_n$ to determine signs (e.g.\ $a_n(2-a_n)\ge 0$).

\begin{itemize}
\item For the invariant-interval step, try $M=$ the larger fixed point of $f$ (solve $L=f(L)$).
\item For monotonicity, the sign of $f(x)-x$ on the invariant interval is the key.
\item Common pitfall: assuming monotonicity from the first few terms without proof by induction.
\end{itemize}

\end{examtip}

\begin{extension}[title={Extension (Tutorial 8, Question 4: $(1+\frac{1}{n})^n$ is increasing)}]

\textbf{Syllabus link:} Chapter 4 --- \S4.3 Monotonic Sequences.

\medskip

\textbf{Proposition.} The sequence $e_n=\left(1+\frac{1}{n}\right)^n$ is increasing and bounded above (by $3$), hence convergent. Its limit is $e$.

\textbf{Proof of monotonicity (calculus approach).}
Define $\phi(x)=x\ln(1+1/x)$ for $x>0$, so that $e_n=e^{\phi(n)}$. Differentiate:
\[
\phi'(x)=\ln\!\left(1+\frac{1}{x}\right)-\frac{1}{x+1}.
\]
Using the inequality $\ln(1+u)\ge \frac{u}{1+u}$ for $u>0$ (with $u=1/x$):
\[
\ln\!\left(1+\frac{1}{x}\right)\ge \frac{1/x}{1+1/x}=\frac{1}{x+1}.
\]
Hence $\phi'(x)\ge 0$, so $\phi$ is increasing, and $e_n=e^{\phi(n)}$ is increasing.

\textbf{Bounded:} By the binomial theorem,
\[
e_n = \sum_{k=0}^{n}\binom{n}{k}\frac{1}{n^k}
\le 1 + \sum_{k=1}^{n}\frac{1}{k!}
\le 1 + \sum_{k=1}^{n}\frac{1}{2^{k-1}}
< 1 + 2 = 3.
\]

By the Monotone Sequence Theorem, $e_n$ converges. The limit is denoted $e\approx 2.71828$.

\end{extension}

%--------------------------------------------------

\subsection{Series}

\paragraph{From a sequence to an ``infinite sum''.}

An infinite series is written $\sum_{n=1}^{\infty} a_n$.

Since we can only add finitely many numbers at a time, the rigorous meaning is via partial sums.

\noindent
\textbf{Why partial sums?}
We know how to add $2$ numbers, $3$ numbers, $\ldots$, $k$ numbers. But ``adding infinitely many numbers'' is not defined by the arithmetic axioms. The only way to give $\sum_{n=1}^{\infty}a_n$ a meaning is to define it as a \emph{limit}: form the finite sums $s_k=a_1+\cdots+a_k$ and ask whether $\{s_k\}$ converges. If it does, the limit \emph{is} the ``infinite sum.'' If it doesn't, the series has no sum.

This is analogous to how we defined the definite integral as a limit of Riemann sums: infinitely many infinitesimal areas added up via a limiting process.

\begin{definition}{Series and partial sums}{def:series}

Given a sequence $\{a_n\}_{n=1}^{\infty}$, define the $k$-th partial sum

\[
s_k=\sum_{n=1}^{k} a_n.
\]

\begin{enumerate}[label=(\roman*), leftmargin=3em]

\item If $\{s_k\}$ converges to a real number $S$, then the series $\sum_{n=1}^{\infty} a_n$ is \emph{convergent} and we write

$\displaystyle \sum_{n=1}^{\infty} a_n=S$.

\item If $\{s_k\}$ diverges, then the series is \emph{divergent}.

\end{enumerate}

\end{definition}

\noindent
\textbf{Key conceptual point.}
The symbol $\sum_{n=1}^{\infty}a_n$ does \textbf{double duty}: it denotes both the formal expression (the ``series'') and, when convergent, its numerical value $S=\lim_{k\to\infty}s_k$.
Always be clear about which meaning you intend.

\begin{examplebox}[title={Worked Example: a divergent series with oscillating partial sums}]

Consider $\displaystyle \sum_{n=1}^{\infty}(-1)^{n-1}=1-1+1-1+\cdots$.

The partial sums are $s_1=1$, $s_2=0$, $s_3=1$, $s_4=0$, \dots, which do not converge.

Hence the series diverges.

\end{examplebox}

\noindent
\textbf{Why the ``$0=1$'' paradox fails (Lecture Example 4.19).}
The ``proof'' that $0=1$ via $0=(1-1)+(1-1)+\cdots = 1+(-1+1)+(-1+1)+\cdots=1$ fails because it \emph{regroups} terms of a divergent series. Regrouping (associativity) is valid for finite sums but not for infinite series unless convergence is established first.

\begin{commonmistake}

Infinite series are \emph{not} finite sums.

Operations such as regrouping terms, cancelling ``$\infty-\infty$'', or treating ``$\cdots$'' algebraically are not justified unless you can translate them into statements about the partial sums $s_k$.

\end{commonmistake}

\begin{theorem}{Geometric series}{thm:geom}

Let $a\ne 0$ and $r\in\mathbb{R}$. The geometric series

\[
\sum_{n=1}^{\infty} a r^{\,n-1}=a+ar+ar^2+\cdots
\]

converges iff $|r|<1$, and in that case

\[
\sum_{n=1}^{\infty} a r^{\,n-1}=\frac{a}{1-r}.
\]

If $|r|\ge 1$, the series diverges.

\end{theorem}

\noindent
\textbf{Why the formula is $\frac{a}{1-r}$.}
The partial sum formula $s_k=\frac{a(1-r^k)}{1-r}$ is derived by the ``multiply-and-subtract'' trick (see Key Formula below). When $|r|<1$, the term $r^k\to 0$ as $k\to\infty$, so $s_k\to \frac{a(1-0)}{1-r}=\frac{a}{1-r}$.

When $|r|\ge 1$, either $r^k$ diverges ($|r|>1$) or $s_k=ka\to\pm\infty$ ($r=1$), or $r^k$ oscillates without converging ($r=-1$).

\begin{keyformula}[title={Geometric partial sum (exam setup)}]

For $r\ne 1$, the $k$-th partial sum is

\[
s_k=a+ar+\cdots+ar^{k-1}
=\frac{a(1-r^k)}{1-r}.
\]
When $|r|<1$, use $r^k\to 0$ to obtain the infinite sum $\frac{a}{1-r}$.

\end{keyformula}

\noindent
\textbf{Derivation of the partial-sum formula.}
Write:
\begin{align*}
s_k &= a + ar + ar^2 + \cdots + ar^{k-1},\\
r\cdot s_k &= ar + ar^2 + ar^3 + \cdots + ar^{k}.
\end{align*}
Subtract: $s_k - rs_k = a - ar^k$, so $s_k(1-r)=a(1-r^k)$, giving $s_k=\frac{a(1-r^k)}{1-r}$.

\begin{examtip}[title={Exam Focus: Identifying geometric series}]

The first step is always to rewrite the general term in the form $a\cdot r^{n-1}$ (or equivalently $\text{const}\cdot r^n$).

\begin{itemize}[leftmargin=*]
\item Identify $a$ (first term) and $r$ (common ratio).
\item Check $|r|<1$ for convergence.
\item If the sum starts at $n=0$ instead of $n=1$: $\sum_{n=0}^{\infty}r^n=\frac{1}{1-r}$.
\item If the sum starts at $n=k$ instead of $n=1$: factor out $r^{k-1}$ or compute as $S_{\text{total}}-s_{k-1}$.
\end{itemize}

\end{examtip}

\begin{examplebox}[title={Worked Example: sum of a geometric series}]

Evaluate $\displaystyle \sum_{n=1}^{\infty}\frac{(-3)^{n-1}}{4^{n}}$.

\emph{Solution.}

Rewrite

\[
\frac{(-3)^{n-1}}{4^{n}}=\frac{1}{4}\left(-\frac{3}{4}\right)^{n-1}.
\]

This is geometric with $a=\frac14$ and $r=-\frac34$ (so $|r|=\frac34<1$).

Therefore,

\[
\boxed{\sum_{n=1}^{\infty}\frac{(-3)^{n-1}}{4^{n}}=\frac{\frac14}{1-(-\frac34)}=\frac{1/4}{7/4}=\frac{1}{7}.}
\]

\end{examplebox}

\begin{examplebox}[title={Worked Example (Lecture Example 4.21): identifying divergent geometric}]

Is $\displaystyle \sum_{n=1}^{\infty}2^{2n}3^{1-n}$ convergent?

\emph{Solution.}
Rewrite: $2^{2n}3^{1-n}=(2^2)^n\cdot 3\cdot 3^{-n}=3\cdot \left(\frac{4}{3}\right)^n=4\cdot\left(\frac{4}{3}\right)^{n-1}$.

This is geometric with $r=\frac{4}{3}>1$, so the series \textbf{diverges}.

\end{examplebox}

\begin{examplebox}[title={Worked Example: telescoping via partial fractions}]

Evaluate $\displaystyle \sum_{n=1}^{\infty}\frac{1}{n(n+1)}$.

\emph{Solution.}

\[
\frac{1}{n(n+1)}=\frac{1}{n}-\frac{1}{n+1}.
\]

Then

\[
s_k=\sum_{n=1}^{k}\left(\frac{1}{n}-\frac{1}{n+1}\right)=1-\frac{1}{k+1}\xrightarrow[k\to\infty]{}1,
\]
so $\displaystyle \boxed{\sum_{n=1}^{\infty}\frac{1}{n(n+1)}=1}$.

\end{examplebox}

\noindent
\textbf{Why telescoping works.}
When partial fractions produce a difference $f(n)-f(n+1)$, consecutive terms in $s_k$ cancel:
\[
s_k = \bigl(f(1)-f(2)\bigr)+\bigl(f(2)-f(3)\bigr)+\cdots+\bigl(f(k)-f(k+1)\bigr)=f(1)-f(k+1).
\]
Only the first and last terms survive---like a collapsing telescope.

\begin{examplebox}[title={Worked Example (Tutorial 8, Q5(iv)): telescoping plus geometric}]

Evaluate $\displaystyle \sum_{n=2}^{\infty}\left(\frac{1}{e^n}+\frac{2}{n^2-1}\right)$.

\emph{Solution.}
Split into two convergent series:

\textbf{Geometric part:}
$\displaystyle\sum_{n=2}^{\infty}\frac{1}{e^n}=\frac{1/e^2}{1-1/e}=\frac{1}{e(e-1)}.$

\textbf{Telescoping part:}
$\frac{2}{n^2-1}=\frac{2}{(n-1)(n+1)}=\frac{1}{n-1}-\frac{1}{n+1}.$

\[
\sum_{n=2}^{N}\left(\frac{1}{n-1}-\frac{1}{n+1}\right)
=\left(1+\frac{1}{2}\right)-\left(\frac{1}{N}+\frac{1}{N+1}\right)\to \frac{3}{2}.
\]

\textbf{Total:}
$\displaystyle\boxed{\sum_{n=2}^{\infty}\left(\frac{1}{e^n}+\frac{2}{n^2-1}\right)=\frac{1}{e(e-1)}+\frac{3}{2}.}$

\end{examplebox}

\begin{examplebox}[title={Worked Example: harmonic series diverges (grouping idea)}]

Show that $\displaystyle \sum_{n=1}^{\infty}\frac{1}{n}$ diverges.

\emph{Solution.}

Let $s_k=\sum_{n=1}^{k}\frac{1}{n}$. For $k=2^m$,

\[
s_{2^m}
=1+\frac12+\left(\frac13+\frac14\right)+\left(\frac15+\cdots+\frac18\right)+\cdots+\left(\frac{1}{2^{m-1}+1}+\cdots+\frac{1}{2^m}\right).
\]

Each block has $2^{j-1}$ terms, each at least $1/2^j$, so each block sum is at least $1/2$.

Thus $s_{2^m}\ge 1+\frac{m}{2}\to\infty$, hence the harmonic series diverges.

\end{examplebox}

\noindent
\textbf{Why the harmonic series is important.}
It is the canonical example showing that $a_n\to 0$ does \emph{not} imply convergence: $\frac{1}{n}\to 0$, yet $\sum \frac{1}{n}=\infty$.
This is why the $n$th-term test (below) can only prove \emph{divergence}, never convergence.

\begin{theorem}{Necessary condition for series convergence}{thm:term-to-zero}

If $\sum_{n=1}^{\infty} a_n$ converges, then $\lim_{n\to\infty} a_n=0$.

\end{theorem}

\noindent
\textbf{Why this is true.}
If $\sum a_n=S$, then $s_k\to S$ and $s_{k-1}\to S$. Since $a_k=s_k-s_{k-1}$, we get $a_k\to S-S=0$.

\begin{theorem}{The $n$th-term test for divergence}{thm:nth-term}

If $\lim_{n\to\infty} a_n$ does not exist or $\lim_{n\to\infty} a_n\ne 0$, then $\sum_{n=1}^{\infty} a_n$ diverges.

\end{theorem}

\noindent
\textbf{Why this is just the contrapositive of Theorem~\ref{thm:term-to-zero}.}
Theorem~\ref{thm:term-to-zero} says: convergent $\Rightarrow$ $a_n\to 0$.
Contrapositive: $a_n\not\to 0$ $\Rightarrow$ divergent.
That's exactly the $n$th-term test.

\begin{commonmistake}[title={Common Mistake: $a_n\to 0$ does NOT imply convergence}]

If $\lim a_n=0$, the $n$th-term test is \textbf{inconclusive}---it tells you nothing.

The harmonic series ($a_n=1/n\to 0$, but diverges) proves this.

You need a stronger test (comparison, ratio, root, integral, etc.)\ to establish convergence when $a_n\to 0$.

\end{commonmistake}

\begin{examtip}

A fast first check for any series $\sum a_n$:

\begin{itemize}

\item Compute $\lim_{n\to\infty} a_n$. If the limit is not $0$ (or does not exist), the series \emph{immediately} diverges by Theorem~\ref{thm:nth-term}.

\item If $\lim a_n=0$, you cannot conclude convergence; you will need additional ideas (covered later in the course).

\end{itemize}

\end{examtip}

\begin{examplebox}[title={Worked Example (Tutorial 8, Q5 collection): quick divergence checks}]

\textbf{(v)} $\displaystyle\sum_{n=1}^{\infty}\frac{n^3+n^2}{n^3-2n+5}$.

$a_n=\frac{n^3+n^2}{n^3-2n+5}\to \frac{1+0}{1-0+0}=1\ne 0$. \textbf{Diverges} by the $n$th-term test.

\textbf{(iii)} $\displaystyle\sum_{n=1}^{\infty}\frac{e^n}{n^2}$.

$a_n=\frac{e^n}{n^2}\to \infty$ (exponential dominates polynomial). \textbf{Diverges} by the $n$th-term test.

\textbf{(vi)} $\displaystyle\sum_{n=1}^{\infty}\left(\frac{3}{5}n+\frac{2}{n}\right)^2$.

$a_n\sim \frac{9}{25}n^2\to\infty$. \textbf{Diverges} by the $n$th-term test.

\end{examplebox}

\begin{theorem}{Linearity of convergent series}{thm:series-linearity}

If $\sum a_n$ and $\sum b_n$ are convergent and $c\in\mathbb{R}$, then:

\begin{enumerate}[label=(\roman*), leftmargin=3em]

\item $\sum (c a_n)$ converges and $\sum (c a_n)=c\sum a_n$.

\item $\sum (a_n+b_n)$ converges and $\sum (a_n+b_n)=\sum a_n+\sum b_n$.

\item $\sum (a_n-b_n)$ converges and $\sum (a_n-b_n)=\sum a_n-\sum b_n$.

\end{enumerate}

\end{theorem}

\noindent
\textbf{Why linearity holds.}
These follow directly from the limit laws for sequences applied to partial sums:
$\sum_{n=1}^{k}(a_n+b_n)=\sum_{n=1}^{k}a_n+\sum_{n=1}^{k}b_n$, and the limit of a sum is the sum of the limits (provided both converge).

\begin{commonmistake}

The linearity rules require that the involved series are already known to be convergent.

It is \emph{not} valid to use these identities to ``prove'' convergence without a separate convergence argument.

\end{commonmistake}

\begin{examplebox}[title={Worked Example (Lecture Example 4.26): combining series}]

Evaluate $\displaystyle \sum_{n=1}^{\infty}\left(\frac{3}{n(n+1)}+\frac{1}{2^n}\right)$.

\emph{Solution.}

Both component series converge:

$\displaystyle\sum_{n=1}^{\infty}\frac{1}{n(n+1)}=1$ (telescoping) and $\displaystyle\sum_{n=1}^{\infty}\frac{1}{2^n}=\frac{1/2}{1-1/2}=1$ (geometric with $r=1/2$).

By linearity:

\[
\sum_{n=1}^{\infty}\left(\frac{3}{n(n+1)}+\frac{1}{2^n}\right)
= 3\cdot 1 + 1 = \boxed{4}.
\]

\end{examplebox}

\begin{examtip}[title={Exam Focus: Finite number of terms does not affect convergence}]

Adding or removing finitely many terms from a series does not change whether it converges or diverges (though it may change the sum).

For example, if $\sum_{n=4}^{\infty}\frac{n}{n^3+1}$ converges, then $\sum_{n=1}^{\infty}\frac{n}{n^3+1}$ also converges (it just adds three extra finite terms).

This is because the partial sums of the two series differ by a fixed constant (the sum of the finitely many added/removed terms), and shifting a sequence by a constant does not affect convergence.

\end{examtip}

\begin{extension}[title={Extension (Beyond Syllabus: $p$-series preview and the integral test connection)}]

\textbf{Syllabus link:} Chapter 4 --- \S4.4 Series.

\medskip

The harmonic series $\sum \frac{1}{n}$ diverges, but what about $\sum \frac{1}{n^2}$ or $\sum \frac{1}{n^p}$ in general?

\textbf{Theorem ($p$-series test).}
The series $\displaystyle\sum_{n=1}^{\infty}\frac{1}{n^p}$ converges if and only if $p>1$.

\textbf{Intuition.}
For $p>1$, the terms $1/n^p$ shrink fast enough that the partial sums stay bounded.
For $p\le 1$, the terms shrink too slowly (at most as fast as $1/n$), and the grouping argument from the harmonic series can be adapted.

This result is typically proved using the integral test (which may appear later in MH1101 or MH2100). The key idea: $\sum \frac{1}{n^p}$ converges iff $\int_1^{\infty}\frac{1}{x^p}\,dx$ converges, and this integral equals $\frac{1}{p-1}$ when $p>1$ and diverges when $p\le 1$.

\end{extension}

%==================================================
% USER COMMENTS (EDITABLE)
%==================================================
% add explanantion to the ftc part especially, avoid giving the definitons directly without any explanaitoin on why are they like that
